\chapter{ಹೆಡರ್ಗಳು}

\section{ಸ್ವಚ್ಛ ವೇದಿಕೆ}
\begin{figure}[h]
\begin{center}
\begin{Scratch}[1]
\scbox{ಅಳಿಸು}{pen}
\end{Scratch}
\end{center}
\caption{ಸ್ಕ್ರಾಚ್ ಪ್ರೋಗ್ರಾಂ ಹೆಡರ್ 1}
\label{hdr_init1}
\end{figure}

\begin{figure}[h]
\begin{center}
\begin{Scratch}[1]
\scbox{ಅಳಿಸು}{pen}
\scbox{\textenglish{x:} \cb[w]{0} \textenglish{y:}\cb[w]{0}ಗೆ ಹೋಗು}{motion}
\end{Scratch}
\end{center}
\caption{ಸ್ಕ್ರಾಚ್ ಪ್ರೋಗ್ರಾಂ ಹೆಡರ್ 2}
\label{hdr_init2}
\end{figure}


\begin{figure}[h]
\begin{center}
\begin{Scratch}[1]
\scbox{ಅಳಿಸು}{pen}
\scbox{\textenglish{x:} \cb[w]{0} \textenglish{y:}\cb[w]{0}ಗೆ ಹೋಗು}{motion}
\scbox{\cb[w]{90} ಯ ದಿಕ್ಕಿನಲ್ಲಿಡು}{motion}
\end{Scratch}
\end{center}
\caption{ಸ್ಕ್ರಾಚ್ ಪ್ರೋಗ್ರಾಂ ಹೆಡರ್ 3}
\label{hdr_init3}
\end{figure}


\begin{figure}[h]
\begin{center}
\begin{Scratch}[1]
\scbox{ಅಳಿಸು}{pen}
\scbox{ಲೇಖನಿಮುಕ್ತ}{pen}
\scbox{\textenglish{x:} \cb[w]{0} \textenglish{y:}\cb[w]{0}ಗೆ ಹೋಗು}{motion}
\scbox{\cb[w]{90} ಯ ದಿಕ್ಕಿನಲ್ಲಿಡು}{motion}
\end{Scratch}
\end{center}
\caption{ಸ್ಕ್ರಾಚ್ ಪ್ರೋಗ್ರಾಂ ಹೆಡರ್ 4}
\label{hdr_init4}
\end{figure}

\section{ಪ್ರೋಗ್ರಾಂ ಉಳಿಸುವುದು}
\begin{figure}[h]
\begin{center}
\includegraphics[scale=2]{Scratchlogo.pdf}
\end{center}
\caption{ಸ್ಕ್ರಾಚ್ ಲೊಗೊ}
\label{hdr_logo}
\end{figure}


\textenglish{File}
ಹೊಸ 
ಯೋಜನೆಯನ್ನು ಉಳಿಸು 
ಬೇಡ 
ಸರಿ 
ರದ್ದುಮಾಡು

\section{ಪ್ರೋಗ್ರಾಂ ಹಂತ ಪ್ರಕ್ರಿಯೆ}

ಚಿತ್ರಸಂಕಲನ

ಪ್ರತಿಯೊಂದು ಹಂತವನ್ನು ಚಲಿಸು 

ಪ್ರತಿಯೊಂದು ಹಂತವನ್ನು ತೋರಿಸು  


\begin{figure}[h]
\begin{center}
\begin{Scratch}[1]
\beginbox{}
\scbox{\cb[w]{100} ಹೆಜ್ಜೆ ಮುಂದೆ ಹೋಗು}{motion}
\turnbox{1}{90}{ಡಿಗ್ರಿಯಷ್ಟು ತಿರುಗು}
\scbox{\cb[w]{100} ಹೆಜ್ಜೆ ಮುಂದೆ ಹೋಗು}{motion}
\turnbox{1}{90}{ಡಿಗ್ರಿಯಷ್ಟು ತಿರುಗು}
\scbox{\cb[w]{100} ಹೆಜ್ಜೆ ಮುಂದೆ ಹೋಗು}{motion}
\turnbox{1}{90}{ಡಿಗ್ರಿಯಷ್ಟು ತಿರುಗು}
\scbox{\cb[w]{100} ಹೆಜ್ಜೆ ಮುಂದೆ ಹೋಗು}{motion}
\end{Scratch}
\end{center}
\caption{ಸ್ಕ್ರಾಚ್ ಪ್ರೋಗ್ರಾಂ ಹಂತ ಪ್ರಕ್ರಿಯೆ}
\label{hdr_sstep}
\end{figure}